% GENERATED FILE. Edit canonical lessons, then run npm run generate:books.
% canonical-source-hash: fnv1a64:902ac737a70df3ad
% canonical-lessons: SA-C411-rathya, SA-C411-apayati, SA-C411-sanketah, SA-C411-rodhah, SA-C411-varjayati

\chapter{Streets and Blockages}
\label{ch:sa-streets-and-blockages}
\hlchaptermodality{411}

\begin{chapteropening}
\textbf{By the end of this chapter:} \emph{I can say how traffic gets round a blockage.}

Complete the last lesson of chapter 411: I can say how traffic gets round a blockage.
\end{chapteropening}

\section[\texorpdfstring{\sk{रथ्या} (rathyā)}{रथ्या (rathyā)}]{\sk{रथ्या} (rathyā) — a carriage road; a street}

\label{lesson:SA-C411-rathya}



\begin{quote}
Before the new one: say the Sanskrit for split open, then the Sanskrit for scrapes.
\end{quote}

\subsection*{You'll want to know: \sk{रथ्या}}
\textbf{\sk{रथ्या}} — \emph{rathyā} — \textquotedblleft{}a carriage road; a street\textquotedblright{}.

\begin{quote}
\sk{रथ्यायां} \sk{वाहनानि} \sk{चलन्ति।}
\end{quote}

\begin{quote}
\emph{rathyāyāṃ vāhanāni calanti.}
\end{quote}

\begin{quote}
Vehicles move along the carriage road.
\end{quote}

A feminine noun.

\begin{cousinweb}
\textbf{\sk{रथ्या}} (\emph{rathyā}) is built on \emph{ratha}, \textquotedblleft{}a chariot\textquotedblright{}, with the ending \emph{-ya}; MW gives \emph{rathyā} as a carriage road, a highway, a street (MW).
\end{cousinweb}

\begin{grammarlens}[title={What you've built}]
One more word for getting around.
\end{grammarlens}

\subsection*{Your turn}
\begin{itemize}
\raggedright
  \item \emph{Say it:} \emph{rathyā}
  \item \emph{Say it:} \emph{rathyā}, once more
  \item \emph{Say it:} \emph{rathyāyāṃ vāhanāni calanti.} — Vehicles move along the carriage road.
  \item \emph{Recall:} say the Sanskrit for split open, then the Sanskrit for scrapes, then say \emph{rathyā} again
\end{itemize}

\begin{tcolorbox}[breakable,colback=teal!4,colframe=teal!35!black,title={Before you move on}]
What does \sk{रथ्या} mean? (\textquotedblleft{}A carriage road; a street\textquotedblright{}.) Say it once more.
\end{tcolorbox}
\section[\texorpdfstring{\sk{अपयाति} (apayāti)}{अपयाति (apayāti)}]{\sk{अपयाति} (apayāti) — moves off; goes away}

\label{lesson:SA-C411-apayati}



\begin{quote}
Before the new one: say the Sanskrit for scrapes, then the Sanskrit for a carriage road.
\end{quote}

\subsection*{You'll want to know: \sk{अपयाति}}
\textbf{\sk{अपयाति}} — \emph{apayāti} — \textquotedblleft{}moves off; goes away\textquotedblright{}.

\begin{quote}
\sk{गजः} \sk{मार्गात्} \sk{अपयाति।}
\end{quote}

\begin{quote}
\emph{gajaḥ mārgāt apayāti.}
\end{quote}

\begin{quote}
The elephant moves off the road.
\end{quote}

Past (it …): \sk{गजः} \sk{मार्गात्} \sk{अपयातवान्}, \emph{gajaḥ mārgāt apayātavān}.

\begin{cousinweb}
\textbf{\sk{अपयाति}} (\emph{apayāti}) is built from \emph{apa}, \textquotedblleft{}away\textquotedblright{}, and the root √yā, \textquotedblleft{}to go\textquotedblright{}; MW gives \emph{apa-yā} as to go away or depart, with the participle \emph{apayāta} (MW).
\end{cousinweb}

\begin{grammarlens}[title={What you've built}]
One more word for getting around.
\end{grammarlens}

\subsection*{Your turn}
\begin{itemize}
\raggedright
  \item \emph{Say it:} \emph{apayāti}
  \item \emph{Say it:} \emph{apayāti}, once more
  \item \emph{Say it:} \emph{gajaḥ mārgāt apayāti.} — The elephant moves off the road.
  \item \emph{Recall:} say the Sanskrit for scrapes, then the Sanskrit for a carriage road, then say \emph{apayāti} again
\end{itemize}

\begin{tcolorbox}[breakable,colback=teal!4,colframe=teal!35!black,title={Before you move on}]
What does \sk{अपयाति} mean? (\textquotedblleft{}Moves off; goes away\textquotedblright{}.) Say it once more.
\end{tcolorbox}
\section[\texorpdfstring{\sk{संकेतः} (saṅketaḥ)}{संकेतः (saṅketaḥ)}]{\sk{संकेतः} (saṅketaḥ) — a signal; a sign}

\label{lesson:SA-C411-sanketah}



\begin{quote}
Before the new one: say the Sanskrit for a carriage road, then the Sanskrit for moves off.
\end{quote}

\subsection*{You'll want to know: \sk{संकेतः}}
\textbf{\sk{संकेतः}} — \emph{saṅketaḥ} — \textquotedblleft{}a signal; a sign\textquotedblright{}.

\begin{quote}
\sk{आरक्षकः} \sk{हस्तेन} \sk{संकेतं} \sk{ददाति।}
\end{quote}

\begin{quote}
\emph{ārakṣakaḥ hastena saṅketaṃ dadāti.}
\end{quote}

\begin{quote}
The policeman gives a signal with his hand.
\end{quote}

A masculine noun.

\begin{cousinweb}
\textbf{\sk{संकेतः}} (\emph{saṅketaḥ}) is built from \emph{sam}, \textquotedblleft{}together\textquotedblright{}, and the root √cit, \textquotedblleft{}to perceive\textquotedblright{}; MW gives \emph{saṃketa} as an agreement or an appointment, and a sign or signal agreed on is that sense (MW).
\end{cousinweb}

\begin{grammarlens}[title={What you've built}]
One more word for getting around.
\end{grammarlens}

\subsection*{Your turn}
\begin{itemize}
\raggedright
  \item \emph{Say it:} \emph{saṅketaḥ}
  \item \emph{Say it:} \emph{saṅketaḥ}, once more
  \item \emph{Say it:} \emph{ārakṣakaḥ hastena saṅketaṃ dadāti.} — The policeman gives a signal with his hand.
  \item \emph{Recall:} say the Sanskrit for a carriage road, then the Sanskrit for moves off, then say \emph{saṅketaḥ} again
\end{itemize}

\begin{tcolorbox}[breakable,colback=teal!4,colframe=teal!35!black,title={Before you move on}]
What does \sk{संकेतः} mean? (\textquotedblleft{}A signal; a sign\textquotedblright{}.) Say it once more.
\end{tcolorbox}
\section[\texorpdfstring{\sk{रोधः} (rodhaḥ)}{रोधः (rodhaḥ)}]{\sk{रोधः} (rodhaḥ) — a blockage; an obstruction}

\label{lesson:SA-C411-rodhah}



\begin{quote}
Before the new one: say the Sanskrit for moves off, then the Sanskrit for a signal.
\end{quote}

\subsection*{You'll want to know: \sk{रोधः}}
\textbf{\sk{रोधः}} — \emph{rodhaḥ} — \textquotedblleft{}a blockage; an obstruction\textquotedblright{}.

\begin{quote}
\sk{मार्गे} \sk{महान्} \sk{रोधः} \sk{अस्ति।}
\end{quote}

\begin{quote}
\emph{mārge mahān rodhaḥ asti.}
\end{quote}

\begin{quote}
There is a big blockage on the road.
\end{quote}

A masculine noun.

\begin{cousinweb}
\textbf{\sk{रोधः}} (\emph{rodhaḥ}) comes from the root √rudh, \textquotedblleft{}to obstruct, to hold back\textquotedblright{}; MW gives \emph{rodha} as stopping, checking, obstructing and blockading (MW).
\end{cousinweb}

\begin{grammarlens}[title={What you've built}]
One more word for getting around.
\end{grammarlens}

\subsection*{Your turn}
\begin{itemize}
\raggedright
  \item \emph{Say it:} \emph{rodhaḥ}
  \item \emph{Say it:} \emph{rodhaḥ}, once more
  \item \emph{Say it:} \emph{mārge mahān rodhaḥ asti.} — There is a big blockage on the road.
  \item \emph{Recall:} say the Sanskrit for moves off, then the Sanskrit for a signal, then say \emph{rodhaḥ} again
\end{itemize}

\begin{tcolorbox}[breakable,colback=teal!4,colframe=teal!35!black,title={Before you move on}]
What does \sk{रोधः} mean? (\textquotedblleft{}A blockage; an obstruction\textquotedblright{}.) Say it once more.
\end{tcolorbox}
\section[\texorpdfstring{\sk{वर्जयति} (varjayati)}{वर्जयति (varjayati)}]{\sk{वर्जयति} (varjayati) — avoids; keeps away from}

\label{lesson:SA-C411-varjayati}



\begin{quote}
Before the new one: say the Sanskrit for a signal, then the Sanskrit for a blockage.
\end{quote}

\subsection*{You'll want to know: \sk{वर्जयति}}
\textbf{\sk{वर्जयति}} — \emph{varjayati} — \textquotedblleft{}avoids; keeps away from\textquotedblright{}.

\begin{quote}
\sk{चालकः} \sk{तं} \sk{मार्गं} \sk{वर्जयति।}
\end{quote}

\begin{quote}
\emph{cālakaḥ taṃ mārgaṃ varjayati.}
\end{quote}

\begin{quote}
The driver keeps away from that road.
\end{quote}

Past (he …): \sk{सः} \sk{तं} \sk{मार्गं} \sk{वर्जितवान्}, \emph{saḥ taṃ mārgaṃ varjitavān}.

\begin{cousinweb}
\textbf{\sk{वर्जयति}} (\emph{varjayati}) is built on the root √vṛj, \textquotedblleft{}to twist, to turn aside\textquotedblright{}, in its causative; MW gives its participle \emph{varjita} as excluded, abandoned, avoided (MW).
\end{cousinweb}

\begin{grammarlens}[title={What you've built}]
One more word for getting around.
\end{grammarlens}

\subsection*{Your turn}
\begin{itemize}
\raggedright
  \item \emph{Say it:} \emph{varjayati}
  \item \emph{Say it:} \emph{varjayati}, once more
  \item \emph{Say it:} \emph{cālakaḥ taṃ mārgaṃ varjayati.} — The driver keeps away from that road.
  \item \emph{Recall:} say the Sanskrit for a signal, then the Sanskrit for a blockage, then say \emph{varjayati} again
\end{itemize}

\begin{tcolorbox}[breakable,colback=teal!4,colframe=teal!35!black,title={Before you move on}]
What does \sk{वर्जयति} mean? (\textquotedblleft{}Avoids; keeps away from\textquotedblright{}.) Say it once more.
\end{tcolorbox}
